As already stated, the forkserver lies inside a \gls{SO} that is injected as dynamic dependence to the binary.
Moreover, as seen in \cref{sec:optimization}, the pointer to the forkserver function is initialized 
(togheter with the \gls{SHM} pointer) in the init function (declared as constructor and executed before the target binary).
During the execution of the ``init'' function the shared object access two environmental variables:
\begin{itemize}
    \item \textit{\_\_AFL\_SHM\_ID} : represent the \gls{SHM} ID setted by the fuzzer (as seen in \cref{sec:state_of_the_art}). 
    \item \textit{\_\_AFL\_PLUSPLUS} : this is used to understand which version of the AFL fuzzer is being used. 
        The tool supports both AFL and AFL++ communication protocol, so this environmental variable is used to choose between the two.
\end{itemize}
Depending on the value of these two environmental variable, 3 different forkserver functions could be used, 
as can be seen in 
\cref{tab:forkserver_modes}
\begin{table}[ht]
    \centering
    \caption{Different execution modes of the forkserver based on environmental variables value}
    \label{tab:forkserver_modes}
    \begin{tabular}{c|c|c}
    \hline
    \textbf{\_\_AFL\_SHM\_ID} & \textbf{\_\_AFL\_PLUSPLUS} & forkserver execution mode\\
    \hline
    NULL & NULL & \textit{standalone mode}\\
    \hline
    NULL & \textit{any value} & \textit{standalone mode}\\
    \hline
    NOT NULL & NULL & \textit{afl++ mode}\\
    \hline
    NOT NULL & 0 & \textit{afl mode}\\
    \hline
    NOT NULL & \textit{any other value} & \textit{afl++ mode}\\
        
    \end{tabular}
\end{table}
\begin{itemize}
    \item \textbf{standalone mode} : this is a modality in which the fuzzer is not present. It consent to run the binary  
    only one time (without the typical loop of the forkserver) without any communication protocol. In this case there is not 
    the \gls{SHM} so the piece of memory that is actually updated is instantiated by the forkserver with a \texttt{malloc} call.\footnote{
        this fact is completely transparent to the target binary which see the relocated memory address as it be a shared memory
    }
    \item \textbf{afl mode} : In this case there is an actual \gls{SHM} and a fuzzer that track the edge coverage.
    \item \textbf{afl++} : Same as the previous point but with a different communication protocol between the forkserver 
        and the fuzzer
\end{itemize}

The \gls{SO} exports two different symbols:
\begin{itemize}
    \item \texttt{uint8\_t *\_\_afl\_area\_ptr} : pointer to the \gls{SHM} (or the local memory instantiated by 
        \texttt{malloc})
    \item \texttt{void (*\_\_forkserver\_start\_ptr)(void)} : pointer to the forkserver function that 
        must be invoked by the main
\end{itemize}
These symbols are resolved inside the target binary in such a way that the binary has access to their value.
Depending on the execution mode, these variables have different values.
\begin{itemize}
    \item \textbf{standalone mode}:\begin{itemize}
        \item \textit{\_\_forkserver\_start\_ptr} : The forkserver function is will return without fork because 
            the target must be executed without tracing. 
        \item \textit{\_\_afl\_area\_ptr} : Allocated through a malloc because it does not need to be a shared map. 
            In this case the edges are traced by the target altought this tracing is not actually used.
            But, for compatibility with the traced case, the \_\_afl\_area\_ptr is kept.
    \end{itemize}
    \item \textbf{afl mode}:\begin{itemize}
        \item \textit{\_\_forkserver\_start\_ptr} : In this case there is an actual communication with the fuzzer. 
            The protocol used can be seen in \cref{alg:afl_forkserver}.
            
        \item \textit{\_\_afl\_area\_ptr} : During the init function to the pointer is associated the \gls{SHM} address 
            through a \texttt{shmat} function, passing as parameter the ID of the \gls{SHM} obtained from 
            the environmental variable.
    \end{itemize}
    \item \textbf{afl++ mode} : This mode works in the same way as the previous but uses a different communication protocol.
        The communication protocol is shown precisely in \cref{fig:aflplusplusforkserver}
        
\end{itemize}

\begin{algorithm}[h]
    \caption{AFL forkserver algorithm}
    \label{alg:afl_forkserver}
    \begin{algorithmic}
        \STATE $\textsc{write}(FORKSRV\_FD+1,0)$
        \WHILE{$\text{true}$}
            \STATE $\textsc{read}(FORKSRV\_FD)$
            \STATE $child\_pid \gets \textsc{fork}()$
            \IF{$child\_pid=0$}
                \STATE $\textsc{close}(FORKSRV\_FD)$ \COMMENT{\small \textcolor{green!60!black}{The child must return from the forkserver and execute the main, so the child process can close the file descriptors}}
                \STATE $\textsc{close}(FORKSRV\_FD+1)$
                \STATE $\text{return}$
            \ELSE
                \STATE $\textsc{write}(FORKSRV\_FD+1,child\_pid)$  \COMMENT{\small \textcolor{green!60!black}{Communicate the PID of the child to the fuzzer}}
                \STATE $\textsc{waitpid}(child\_pid)$
            \ENDIF
        \ENDWHILE
    \end{algorithmic}
\end{algorithm}

\textbf{libforkserver.so} has also a destructor function called \texttt{void afl\_forkserver\_cleanup(void)}.
A function declared as destructor is a function with the attribute \texttt{\_\_attribute\_\_((destructor))}.
A reference to a function declared with the previous attribute is placed inside the \textit{.fini\_array} and is executed 
after the main.
This function is used to detatch the \gls{SHM} from the current process (in case of afl/afl++ execution mode), or to 
free the memory in case of standalone execution mode.